\documentclass[12pt, oneside, a4paper, chapter=TITLE,section=TITLE]{abntex2}

\usepackage{abntex2-ifrj} % pacote com customizações da classe abntex2 para uso no IFRJ-CVR

%Carregando alguns pacotes extras
\usepackage{float}

% ---
% Informações de dados para CAPA e FOLHA DE ROSTO
% ---
\titulo{Modelo de Projeto de pesquisa para submissão ao Comitê de Apoio à Pesquisa do \textit{campus} Volta Redonda}
\tipotrabalho{desenvolvimento tecnológico e inovação} % colocar "pesquisa" ou "desenvolvimento tecnológico e inovação"
\autor{Nome do(a) coordenador(a) do projeto}
\local{Volta Redonda}
\data{\today} % colocar a data de entrega ao CAP
\programa{Comitê de Apoio à Pesquisa}
% Preâmbulo 
\preambulo{Projeto de \imprimirtipotrabalho\  apresentado ao \imprimirprograma do \textit{Campus} Volta Redonda como requisito para validação institucional de seu desenvolvimento no \textit{campus}.
} % O preambulo deve conter o tipo do trabalho, o objetivo, o nome da instituição e a área de concentração

\linhapesquisa{Otimização topológica} % colocar a linha de pesquisa à qual o projeto está vinculado, conforme o grupo de pesquisa do qual o pesquisador coordenador é integrante.

% --------------------------------
\makeindex % compila o índice remissivo (opcional)
% --------------------------------

% --------------------------------
% Início do documento
% --------------------------------
\begin{document}

% Retira espaço extra obsoleto entre as frases.
\frenchspacing

% --------------------------------
% ELEMENTOS PRÉ-TEXTUAIS
% --------------------------------
\pretextual
% ---
\imprimircapa % Capa
\imprimirfolhaderosto % Folha de rosto
% ---

% ---
% NOTA DA ABNT NBR 15287:2011, p. 4:
%  ``Se exigido pela entidade, apresentar os dados curriculares do autor em folha ou página distinta após a folha de rosto.''
% ---

% ---
% lista de ilustrações - Comente o código abaixo caso não tenha nenhuma imagem 
\pdfbookmark[0]{\listfigurename}{lof} \listoffigures*
\cleardoublepage
% ---

% ---
% lista de tabelas - Comente o código abaixo caso não tenha nenhuma tabela 
\pdfbookmark[0]{\listtablename}{lot}\listoftables*
\cleardoublepage
% ---

% ---
% inserir lista de quadros - Comente o código abaixo caso não tenha nenhum quadro 
\pdfbookmark[0]{\listofquadrosname}{loq} \listofquadros*
\cleardoublepage
% ---

% ---
% inserir lista de abreviaturas e siglas (opcional)
\begin{siglas}
  \item[ABNT] Associação Brasileira de Normas Técnicas
  \item[abnTeX] ABsurdas Normas para TeX
\end{siglas}
% ---


% ---
% inserir lista de símbolos
\begin{simbolos}
  \item[$ \Gamma $] Letra grega Gama
  \item[$ \Lambda $] Lambda
  \item[$ \zeta $] Letra grega minúscula zeta
  \item[$ \in $] Pertence
\end{simbolos}
% ---



% ---
% sumario
\pdfbookmark[0]{\contentsname}{toc}
\tableofcontents*
% \cleardoublepage
% ---


% --------------------------------
% ELEMENTOS TEXTUAIS
\textual
% --------------------------------

\chapter[Introdução]{Introdução}
A introdução é a parte inicial do texto, tendo como propósito apresentar o tema, o(s) problema(s) e a(s) hipótese(s). Deve ser clara e objetiva, expondo, de forma sucinta, a natureza e a intencionalidade do trabalho. 

Este modelo foi elaborado tendo como base a classe \textsf{abntex2} e do pacote \textsf{abntex2cite} visando atender os requisitos da ABNT NBR 15287:2011 \emph{Informação e documentação - Projeto de
pesquisa - Apresentação}.

Cabe aos pesquisadores observarem o manual da classe \textsf{abntex2} e as normas pertinentes. Uma lista das normas observadas pelo \abnTeX\ é apresentada em \citeonline{abntex2classe}.

\lipsum[1-5]


\chapter{Objetivos}
Devem-se delimitar o(s) objetivo(s) geral(is) e os objetivos específicos da pesquisa.


\chapter{Justificativa}
Apresentam-se, nessa parte, as questões que delimitam a escolha do tema, o contexto relativo à questão investigada, sua origem e relevância quanto a aspectos científicos, pedagógicos e/ou sociais.

\chapter{Fundamentação Teórica}
São apresentados e discutidos os conceitos relacionados à questão pesquisada. Também podem ser avaliadas criticamente obras e pesquisas realizadas sobre o tema.

\section{Exemplo de seção}
\subsection{exemplo de subseção}


\chapter{Metodologia}
Devem-se apresentar o tipo de pesquisa, os instrumentos de investigação, a análise de dados, as limitações do estudo etc.

Nessa parte, explica-se como serão realizados os objetivos propostos, com base nas hipóteses levantadas. 


\chapter{Plano de Trabalho}
Deve-se descrever as etapas do trabalho com os respectivos prazos previstos.

Recomenda-se que seja feita a descrição das etapas e o coronograma de execução conforme exemplo abaixo.



\begin{etapas}
	\item \label{e1} Descrição da etapa 1
	\item \label{e2} Descrição da etapa 2
	\item \label{e3} Sed consequat tellus et tortor. Ut tempor laoreet quam. Nullam id wisi a liberotristique  semper.  Nullam  nisl  massa,  rutrum  ut,  egestas  semper,  mollis  id,  leo.  Nullaac  massa  eu  risus  blandit  mattis.  Mauris  ut  nunc.  In  hac  habitasse  platea  dictumst.Aliquam eget tortor. Quisque dapibus pede in erat.
	\item \label{e4} Sed consequat tellus et tortor. Ut tempor laoreet quam. 
	\item \label{e5} Sed consequat tellus et tortor. Ut tempor laoreet quam. 
	\item \label{e6} Sed consequat tellus et tortor. Ut tempor laoreet quam. 
	\item \label{e7} Sed consequat tellus et tortor. Ut tempor laoreet quam. 
\end{etapas}

\begin{center}

\begin{ganttchart}[vgrid,hgrid]{1}{24}
\gantttitle{2021}{12} %12 meses em 2021
\gantttitle{2022}{12} \\ %12 meses em 2022
    \gantttitlelist{1,...,12}{1} % meses 1 ao 12 de 2021
    \gantttitlelist{1,...,12}{1} \\ % meses 1 ao 12 de 2022
        \ganttbar{\ref{e1}}{1}{7} \\
        \ganttbar{\ref{e2}}{7}{12} \\
        \ganttbar{\ref{e3}}{13}{18} \\
        \ganttbar{\ref{e4}}{18}{24}
\end{ganttchart}

\begin{ganttchart}[vgrid,hgrid]{1}{12}
\gantttitle{2023}{12}\\ %12 meses em 2023
    \gantttitlelist{1,...,12}{1} \\ % meses 1 ao 12 de 2023
        \ganttbar{\ref{e5}}{1}{7} \\
        \ganttbar{\ref{e6}}{8}{10} \\
        \ganttbar[progress=65]{\ref{e7}}{10}{12}
\end{ganttchart}

\end{center}

\chapter{Resultados Esperados}
Descrevem-se as expectativas de resultados ao final do trabalho.


% ---------------------------------
% ELEMENTOS PÓS-TEXTUAIS
% ---------------------------------
\phantompart
\postextual

% ---------------------------------
% Referências bibliográficas
% ---------------------------------
\bibliography{references}

% ---------------------------------
% Glossário
% ---------------------------------
%
% Consulte o manual da classe abntex2 para orientações sobre o glossário.
%
% \glossary

% ---------------------------------
% Apêndices  Elemento opcional
% ---------------------------------

\begin{apendicesenv} % Inicia os apêndices

\partapendices % Imprime uma página indicando o início dos apêndices

% ---------------------------------
\chapter{Quisque libero justo}
% ---------------------------------

\lipsum[50]

% ---------------------------------
\chapter{Nullam elementum urna vel imperdiet sodales elit ipsum pharetra ligula ac pretium ante justo a nulla curabitur tristique arcu eu metus}
% ---------------------------------
\lipsum[55-57]

\end{apendicesenv} % finaliza os apêndices
% ---


% ----------------------------------------------
% Anexos - Elemento opcional
% ----------------------------------------------
\begin{anexosenv} % Inicia os anexos

\partanexos % Imprime uma página indicando o início dos anexos

% ---
\chapter{Morbi ultrices rutrum lorem.}
% ---
Pode-se anexar, ao final do projeto, material suplementar de sustentação à pesquisa: questionários, fotocópias, fotografias etc.

\lipsum[30]

% ---
\chapter{Cras non urna sed feugiat cum sociis natoque penatibus et magnis dis parturient montes nascetur ridiculus mus}
% ---

\lipsum[31]

\end{anexosenv} % fim dos anexos

%------------------------------------------------------
% INDICE REMISSIVO - elemento opcional
%------------------------------------------------------
\phantompart
\printindex

\end{document}
